\section{From a common shift to real material perturbations}
\label{sec:channels}

\subsection{Projection and operator decomposition}

Let $P_c$ project onto a localized low-energy or correlated subspace, such as
transition-metal $d$ orbitals and ligand $p$ orbitals.  The projected DFPT
perturbation is
\begin{equation}
  D_{\nu\vq}=P_c\,\partial_{\nu\vq}V_{\ks}\,P_c.
  \label{eq:projectedD}
\end{equation}
For each correlated site $i$, split its on-site block into an orbital trace and a
traceless remainder,
\begin{align}
  D^{\mathrm{charge}}_{i}
  &= \frac{\Tr_{\mathrm{orb}}D_{ii}}{N_{\mathrm{orb}}}\,I_i,
  \label{eq:chargechannel}\\
  D^{\mathrm{internal}}_{i}
  &= D_{ii}-D^{\mathrm{charge}}_i.
  \label{eq:internalchannel}
\end{align}
The off-site blocks $D_{ij}$, $i\ne j$, describe changes in hopping and
hybridization.  If needed, extra channels represent derivatives of interaction
parameters, such as $\partial U/\partial u$ and $\partial J/\partial u$.  Thus
\begin{equation}
  D_{\nu\vq}
  =D^{\mathrm{charge}}_{\nu\vq}
  +D^{\mathrm{internal}}_{\nu\vq}
  +D^{\mathrm{nonlocal}}_{\nu\vq}
  +D^{\mathrm{interaction}}_{\nu\vq}.
  \label{eq:channeldecomp}
\end{equation}
This decomposition is basis-dependent, so the project must state the Wannier or
projector construction and test robustness against a reasonable basis change.

\subsection{Why a local common operator does not shift all bands equally}

Suppose a phonon produces the same local potential $D$ on all orbitals of the
correlated subspace.  A Bloch state still shifts by
\begin{equation}
  \delta\varepsilon_{n\vk}
  =D\langle\psi_{n\vk}|P_c|\psi_{n\vk}\rangle,
  \label{eq:bandweight}
\end{equation}
which depends on its correlated-orbital weight.  At finite momentum,
\begin{equation}
  g_{mn}(\vk,\vq)
  =D\langle\psi_{m,\vk+\vq}|P_c|\psi_{n\vk}\rangle,
  \label{eq:finiteqweight}
\end{equation}
which also depends on phases and interband overlaps.  DFPT already computes these
geometric and hybridization effects.  The proposed correction concerns the
many-body renormalization of the operator, not a replacement of all matrix
elements by a common number.

\subsection{Evidence from correlated materials}

DFT+DMFT calculations provide a useful controlled contrast because they evaluate
the derivative of a local, frequency-dependent correlated self-energy
\citep{Abramovitch2025}.  For SrVO$_3$ with $U=4.5$ eV and $J=0.675$ eV, the
reported zero-frequency coupling for an M-point Jahn--Teller mode increases from
about $44$ meV in DFPT to $87$ meV in DFT+DMFT.  The corresponding R-point
breathing-mode coupling changes from about $58$ meV to $50$ meV.  The same
correlated material therefore contains a strongly corrected orbital-splitting
mode and a comparatively stable charge-like breathing mode.

For hole-doped CaCuO$_2$ at a representative interaction $U=3.1$ eV, the reported
zero-frequency half- and full-breathing couplings change from approximately
$70$ to $76$ meV and from $53$ to $45$ meV, respectively
\citep{Abramovitch2025}.  These static changes are modest even though the
quasiparticle weight is reduced.  At larger $U=4.7$ eV, however, the same study
finds strong vertex frequency dependence on the phonon energy scale.  Agreement at
$\omega=0$ is then insufficient to validate a dynamic scattering calculation.

These cases support a mode-resolved interpretation, but they do not yet prove the
specific decomposition in \cref{eq:channeldecomp}.  Two further methods establish
that standard DFPT has genuine, computable limitations:

\begin{itemize}
  \item DFPT+$U$ corrects both the reference electronic structure and the response
  of Hubbard occupations.  In CoO, this changes the polar Fr\"ohlich interaction
  relative to ordinary DFPT \citep{Zhou2021}.
  \item GW perturbation theory differentiates a nonlocal, energy-dependent GW
  self-energy.  It produces correlation-enhanced electron--phonon interactions in
  Ba$_{1-x}$K$_x$BiO$_3$ \citep{Li2019}.
\end{itemize}

These are not competing anecdotes.  They identify distinct missing channels:
local correlated occupations in DFPT+$U$, local dynamic self-energy in DFT+DMFT,
and nonlocal quasiparticle self-energy in GW perturbation theory.
Earlier correlated-electron analyses likewise show that vertex corrections can be
strongly momentum- and channel-dependent \citep{Grilli1994}.

\subsection{What ``DFPT works in a strongly correlated material'' can mean}

A material may be strongly correlated according to its spectral weight, Hubbard
bands, or mass enhancement while a particular phonon derivative
$\partial_u\Sigma^{\mathrm{corr}}$ remains small or nearly common within the
relevant low-energy subspace.  Conversely, a seemingly small traceless orbital
component can be strongly amplified near an orbital instability.  The correction
is controlled by the product of perturbation weight and channel susceptibility,
not by either factor alone.

Accordingly, a successful DFPT result in one mode should not be generalized to
all modes of that material.  The correct unit of comparison is a
material--phonon--momentum--frequency tuple.
